%-----------------------------------------------------------------------------------------------------------------------------------------------%
%	The MIT License (MIT)
%
%	Copyright (c) 2021 Jitin Nair
%
%	Permission is hereby granted, free of charge, to any person obtaining a copy
%	of this software and associated documentation files (the "Software"), to deal
%	in the Software without restriction, including without limitation the rights
%	to use, copy, modify, merge, publish, distribute, sublicense, and/or sell
%	copies of the Software, and to permit persons to whom the Software is
%	furnished to do so, subject to the following conditions:
%	
%	THE SOFTWARE IS PROVIDED "AS IS", WITHOUT WARRANTY OF ANY KIND, EXPRESS OR
%	IMPLIED, INCLUDING BUT NOT LIMITED TO THE WARRANTIES OF MERCHANTABILITY,
%	FITNESS FOR A PARTICULAR PURPOSE AND NONINFRINGEMENT. IN NO EVENT SHALL THE
%	AUTHORS OR COPYRIGHT HOLDERS BE LIABLE FOR ANY CLAIM, DAMAGES OR OTHER
%	LIABILITY, WHETHER IN AN ACTION OF CONTRACT, TORT OR OTHERWISE, ARISING FROM,
%	OUT OF OR IN CONNECTION WITH THE SOFTWARE OR THE USE OR OTHER DEALINGS IN
%	THE SOFTWARE.
%	
%
%-----------------------------------------------------------------------------------------------------------------------------------------------%

%----------------------------------------------------------------------------------------
%	DOCUMENT DEFINITION
%----------------------------------------------------------------------------------------

% article class because we want to fully customize the page and not use a cv template
\documentclass[a4paper,11pt]{article}

%----------------------------------------------------------------------------------------
%	PACKAGES
%----------------------------------------------------------------------------------------
\usepackage{url}
\usepackage{parskip}

\RequirePackage{color}
\RequirePackage{graphicx}
\usepackage[usenames,dvipsnames]{xcolor}
\usepackage[top=1.1cm, bottom=1.1cm, left=1.2cm, right=1.2cm]{geometry}

\usepackage{tabularx}
\usepackage{enumitem}

\newcolumntype{C}{>{\centering\arraybackslash}X}

\usepackage{supertabular}
\newlength{\fullcollw}
\setlength{\fullcollw}{0.47\textwidth}

\usepackage{titlesec}
\usepackage{multicol}
\usepackage{multirow}

\titleformat{\section}{\large\scshape\raggedright}{}{0em}{}[\titlerule]
\titlespacing{\section}{0pt}{5pt}{5pt}

\usepackage[french]{babel}
\usepackage{csquotes}
\usepackage[style=authoryear,sorting=ynt, maxbibnames=2]{biblatex}

\usepackage[unicode, draft=false]{hyperref}
\definecolor{linkcolour}{rgb}{0,0.2,0.6}
\hypersetup{colorlinks,breaklinks,urlcolor=linkcolour,linkcolor=linkcolour}
\addbibresource{citations.bib}
\setlength\bibitemsep{0.4em}

\usepackage{fontawesome5}

\setlength{\parskip}{0pt}

%----------------------------------------------------------------------------------------
%	BEGIN DOCUMENT
%----------------------------------------------------------------------------------------
\begin{document}
\pagestyle{empty}

%----------------------------------------------------------------------------------------
%	TITLE
%----------------------------------------------------------------------------------------
\begin{tabularx}{\linewidth}{@{} C @{}}
\Huge{Pierre BELAMRI} \\[2pt]
\small{
\href{https://github.com/Rxzh}{\raisebox{-0.05\height}\faGithub\ Rxzh} \ $|$ \
\href{https://linkedin.com/in/pierre-r-belamri-4a8484191}{\raisebox{-0.05\height}\faLinkedin\ Pierre Belamri} \ $|$ \
\href{https://pbelamri.com}{\raisebox{-0.05\height}\faGlobe\ pbelamri.com} \ $|$ \
\href{https://scholar.google.com/citations?user=8xd25vsAAAAJ&hl=en&oi=ao}{\raisebox{-0.05\height}\faGraduationCap\ Google Scholar} \ $|$ \
\href{mailto:pierre.belamri@protonmail.com}{\raisebox{-0.05\height}\faEnvelope\ pierre.belamri@protonmail.com} \ $|$ \
\href{tel:+33750934949}{\raisebox{-0.05\height}\faMobile\ +33.7.50.93.49.49}
}
\end{tabularx}

%----------------------------------------------------------------------------------------
%	SUMMARY
%----------------------------------------------------------------------------------------
\vspace{-2pt}
\section{Profil}
\vspace{1pt}
{\small Doctorant (soutenance en 2026) à Mines Paris -- PSL / CNRS, spécialisé en \textbf{self-supervised et multimodal machine learning} appliqué aux données scientifiques. Mes travaux portent sur l'apprentissage de représentations latentes à partir de données de microstructure et d'essais mécaniques (EBSD, nanoindentation, émission acoustique) pour l'inférence de l'élastoplasticité à l'échelle sub-granulaire. Chercheur invité à UIUC et UCSB. Chargé de cours sur les Transformers à Mines Paris. Fondateur d'une société SaaS d'IA (clients Fortune 500).}

%----------------------------------------------------------------------------------------
%	EXPERIENCE
%----------------------------------------------------------------------------------------
\section{Expérience professionnelle}

\begin{tabularx}{\linewidth}{ @{}l r@{} }
\textbf{Étudiant doctorant --- CNRS \& Mines Paris (PSL), Paris} & \hfill janv. 2024 -- aujourd'hui \\[2pt]
\multicolumn{2}{@{}X@{}}{\small
\begin{minipage}[t]{\linewidth}
    \begin{itemize}[nosep,after=\strut, leftmargin=1em, itemsep=1.5pt]
        \item[--] \textbf{Thèse :} \textit{Machine learning pour la détection d'anomalies en nanoindentation afin d'identifier l'élastoplasticité cristalline à l'échelle sub-granulaire.} Encadrants : D. Ryckelynck, D. Texier.
        \item[--] Développement d'un pipeline self-supervised pour apprendre des représentations latentes structurées de données de microstructure en haute dimension (cartes d'orientation EBSD, courbes P-h, signaux acoustiques).
        \item[--] Conception d'un pipeline de recalage multimodal couplant EBSD, nanoindentation et émission acoustique en un jeu de données unifié (titane enrichi en oxygène) ; article en préparation.
        \item[--] Pré-entraînement d'un Quaternion Vision Transformer sur 2M de patchs EBSD synthétiques (128$\times$128$\times$4 px) ; \textbf{premier auteur}, \textit{Materials \& Design} (2025). Présenté à \textbf{TMS AIM2025}.
        \item[--] HPC Jean-Zay (IDRIS) ; plus gros entraînements : 8$\times$A100 et 4$\times$H100. Code et jeux de données entièrement open source.
    \end{itemize}
    \end{minipage}
    }
\end{tabularx}

\begin{tabularx}{\linewidth}{ @{}l r@{} }
\textbf{Étudiant doctorant invité --- University of Illinois Urbana-Champaign (UIUC)} & \hfill oct. 2025 -- nov. 2025 \\[2pt]
\multicolumn{2}{@{}X@{}}{\small
\begin{minipage}[t]{\linewidth}
    \begin{itemize}[nosep,after=\strut, leftmargin=1em, itemsep=1.5pt]
        \item[--] Séjour de recherche dans le groupe du Prof.\ J.-C.\ Stinville (Materials Science \& Engineering).
        \item[--] \textbf{2\textsuperscript{e} place / 77 équipes internationales} + \textbf{Polaron Sponsor Prize}, 2025 AI/ML for Microscopy Hackathon : développement d'un pipeline ML d'acquisition EBSD automatisée, orientant l'échantillonnage spatial vers les régions out-of-distribution via un prior sur l'espace latent basé sur la distance de Mahalanobis, réduisant fortement l'exposition au faisceau d'électrons. \href{https://kaliningroup.github.io/mic_hackathon_2/submissions/}{Soumission}
    \end{itemize}
    \end{minipage}
    }
\end{tabularx}

\begin{tabularx}{\linewidth}{ @{}l r@{} }
\textbf{Étudiant doctorant invité --- UC Santa Barbara (UCSB)} & \hfill mai 2025 -- juin 2025 \\[2pt]
\multicolumn{2}{@{}X@{}}{\small
\begin{minipage}[t]{\linewidth}
    \begin{itemize}[nosep,after=\strut, leftmargin=1em, itemsep=1.5pt]
        \item[--] Séjour de recherche dans le groupe du Prof.\ Sam Daly (Mechanical Engineering).
        \item[--] Intégration de l'émission acoustique dans des dispositifs d'essais micro-mécaniques ; extraction de features latentes à partir des signaux d'EA pour des architectures ML multimodales.
        \item[--] Coordination stratégique d'une initiative commune entre Mines Paris, UCSB, la NASA et le Synchrotron SOLEIL sur la caractérisation avancée des matériaux.
    \end{itemize}
    \end{minipage}
    }
\end{tabularx}

\begin{tabularx}{\linewidth}{ @{}l r@{} }
\textbf{Chargé de cours --- Mines Paris -- PSL, Paris} & \hfill 2024 -- aujourd'hui \\[2pt]
\multicolumn{2}{@{}X@{}}{\small Enseignement du cours \textbf{Transformers} aux élèves ingénieurs (niveau MSc) : architecture, mécanismes d'attention, pré-entraînement et fine-tuning.}
\end{tabularx}

\begin{tabularx}{\linewidth}{ @{}l r@{} }
\textbf{Fondateur --- VH Systems / EngageAI, Paris} & \hfill mai 2024 -- aujourd'hui \\[2pt]
\multicolumn{2}{@{}X@{}}{\small Conception et commercialisation d'\textbf{EngageAI}, une plateforme d'intelligence OSINT sur les parties prenantes, basée sur l'IA (agents de recherche web, GCP, ReactJS). 4 clients actifs dont 3 entreprises du Fortune 500 (sciences de la vie, pharma, télécommunications). Entièrement autofinancée.}
\end{tabularx}

\begin{tabularx}{\linewidth}{ @{}l r@{} }
\textbf{Consultant en Machine Learning --- Université de Zurich} & \hfill avr. 2024 -- juin 2025 \\[2pt]
\multicolumn{2}{@{}X@{}}{\small
\begin{minipage}[t]{\linewidth}
    \begin{itemize}[nosep,after=\strut, leftmargin=1em, itemsep=1.5pt]
        \item[--] Fine-tuning d'une tête de représentation sur \textbf{NV-Embed-v2} à partir d'un corpus constitué sur l'élection présidentielle américaine de 2024 ; l'espace latent s'organise spontanément selon le compas politique, sans aucune étiquette a priori.
        \item[--] Mise en place d'une infrastructure d'entraînement multi-GPU ; conseil sur le réglage des hyperparamètres et le prétraitement des données auprès de chercheurs en science politique.
    \end{itemize}
    \end{minipage}
    }
\end{tabularx}

\begin{tabularx}{\linewidth}{ @{}l r@{} }
\textbf{Ingénieur ML --- Consultys, Lausanne} & \hfill mars 2023 -- nov. 2023 \\[2pt]
\multicolumn{2}{@{}X@{}}{\small Déploiement de LLMs open source sur les instances cloud des clients (Docker). Développement d'un agent conversationnel multitâche et modulaire pour des Research Scientists. Mémoire de master sur la \textbf{quantization de LLMs} (GPTQ et variantes).}
\end{tabularx}

\begin{tabularx}{\linewidth}{ @{}l r@{} }
\textbf{Stagiaire ML (BI) --- Amazon, Berlin} & \hfill mars 2022 -- juil. 2022 \\[2pt]
\multicolumn{2}{@{}X@{}}{\small Entraînement d'un classifieur NLP d'offres d'emploi concurrentes ; déploiement de bout en bout sur AWS (CDK, Docker, ETL) avec une inférence ARM optimisée en coût.}
\end{tabularx}

\begin{tabularx}{\linewidth}{ @{}l r@{} }
\textbf{Ingénieur logiciel stagiaire --- PMI, Lausanne} & \hfill sept. 2021 -- févr. 2022 \\[2pt]
\multicolumn{2}{@{}X@{}}{\small Scrapers Selenium sur Tor pour la surveillance du commerce illicite ; analyse de graphes et renseignement sur le dark web pour la lutte contre la contrebande de tabac.}
\end{tabularx}

\begin{tabularx}{\linewidth}{ @{}l r@{} }
\textbf{Stagiaire de recherche --- CAOR Mines Paris, Paris} & \hfill mars 2021 -- avr. 2021 \\[2pt]
\multicolumn{2}{@{}X@{}}{\small Computer vision + Deep RL pour la manipulation par bras robotique : \textbf{98,3\,\% de réussite} en moins de 30k epochs d'entraînement. \href{https://arxiv.org/pdf/2105.09719.pdf}{Preprint.}}
\end{tabularx}

%----------------------------------------------------------------------------------------
%	PROJECTS
%----------------------------------------------------------------------------------------
\section{Projets}

\begin{tabularx}{\linewidth}{ @{}l r@{} }
\textbf{CRAFT-JEPA} (recherche indépendante) & \hfill \href{https://github.com/Rxzh/craft-jepa}{GitHub} \\[2pt]
\multicolumn{2}{@{}X@{}}{\small Fine-tuning de \textbf{V-JEPA} sur le dataset VPT (4,5 To de parties de Minecraft + commandes) pour étudier la planification à long horizon en environnement stochastique. Premier entraînement convergé ; évaluation de l'architecture en cours. Entraîné sur 4$\times$H100 (autofinancé).}
\end{tabularx}

%----------------------------------------------------------------------------------------
%	EDUCATION
%----------------------------------------------------------------------------------------
\section{Formation}
\begin{tabularx}{\linewidth}{@{}l X@{}}
2024 -- 2026 & \textbf{Doctorat (PhD)} en Machine Learning appliqué à la science des matériaux --- \textbf{Mines Paris -- PSL} / CNRS \\[-2pt]
2019 -- 2023 & \textbf{Diplôme d'Ingénieur \& MSc} (Ingénieur civil -- Executive Engineering) --- \textbf{Mines Paris -- PSL} \\[-2pt]
2017 -- 2019 & Classes préparatoires --- \textbf{Lycée Marcel Sembat} \\
\end{tabularx}

%----------------------------------------------------------------------------------------
%	SKILLS
%----------------------------------------------------------------------------------------
\section{Compétences}
\begin{tabularx}{\linewidth}{@{}l X@{}}
Recherche       & \small{Self-supervised learning, multimodal representation learning, détection d'anomalies, transformers, CNNs, RL, LLM quantization, NLP} \\[-2pt]
ML Stack        & \small{PyTorch, TensorFlow, HuggingFace, scikit-learn $|$ SLURM/HPC (Jean-Zay), entraînement distribué (8$\times$A100, 4$\times$H100), AWS, GCP, Docker} \\[-2pt]
Ingénierie      & \small{Python, JavaScript/ReactJS, SQL, C\#/Unity} \\[-2pt]
Langues         & \small{Français (langue maternelle), anglais (courant)} \\
\end{tabularx}

%----------------------------------------------------------------------------------------
%	ACHIEVEMENTS
%----------------------------------------------------------------------------------------
\section{Distinctions}
\begin{tabularx}{\linewidth}{@{}l X@{}}
Gaming compétitif & \small{Classé \textbf{14\textsuperscript{e} mondial} en Minecraft PvP compétitif (Kohi, 2016) ; membre clé de \textit{VorHerr}, alors équipe PvP n\textsuperscript{o}\,1 mondiale.} \\
\end{tabularx}

%----------------------------------------------------------------------------------------
%	PUBLICATIONS
%----------------------------------------------------------------------------------------
\section{Publications}
\begin{refsection}[citations.bib]
\nocite{*}
\printbibliography[heading=none]
\end{refsection}

\vspace{2pt}
\center{\footnotesize Dernière mise à jour : \today}

\end{document}